\chapter{Estándares, Prácticas, Convenciones y Métricas}
\label{sec:estandares}

\section{Estándares de Código}
\label{sec:estandares-codigo}

El proyecto SpeechNotes sigue los siguientes estándares y convenciones:

\begin{longtable}{lp{4cm}p{6cm}}
\caption{Estándares de código aplicados}
\label{tab:estandares-codigo}\\
\toprule
\textbf{Lenguaje} & \textbf{Estándar} & \textbf{Herramienta de verificación}\\
\midrule
\endfirsthead
\toprule
\textbf{Lenguaje} & \textbf{Estándar} & \textbf{Herramienta de verificación}\\
\midrule
\endhead
\bottomrule
\endfoot

Python & PEP 8 & Ruff / SonarCloud\\

TypeScript & ESLint + Prettier & SonarCloud\\

Arquitectura & Patrones GoF (Singleton, Factory, Adapter, Facade, Strategy, State, Observer) & Revisiones de código\\

Base de datos & Prisma ORM + MongoDB & Linting de esquema\\

\end{longtable}

\section{Métricas SQA}
\label{sec:metricas}

Las siguientes métricas fueron recolectadas para evaluar la calidad del producto y del proceso de pruebas:

\subsection{Porcentaje de casos pasados vs.\ fallados}
\label{subsec:porcentaje-casos}

\[
\%\text{Pasados} = \frac{\text{Casos pasados}}{\text{Total casos}} \times 100
\]

\subsection{Densidad de defectos}
\label{subsec:densidad-defectos}

\[
\text{Densidad de defectos} = \frac{\text{Número de defectos encontrados}}{\text{Número de casos de prueba ejecutados}}
\]

\subsection{Cobertura de requisitos}
\label{subsec:cobertura-rf}

\[
\%\text{Cobertura RF} = \frac{\text{Requisitos con al menos un caso asociado}}{\text{Total requisitos funcionales}} \times 100
\]

\subsection{Periodicidad de reporte}
\label{subsec:periodicidad}

Las métricas fueron reportadas al finalizar el sprint de calidad (21 de julio de 2026), consolidadas en el Apéndice~E (Informe de Cierre).
